\section{Mathematical Foundation}
\label{sec:mathematical-foundation}

\subsection{Axiomatic Foundation}
\label{subsec:axiomatic-foundation}

The QR theory rests on eight fundamental axioms that establish the relationship between the higher-dimensional information space $\mathcal{I}$ and the observable four-dimensional spacetime $M^{(4)}$. These axioms provide the mathematical foundation for all subsequent developments and ensure internal consistency of the theoretical framework.

\begin{qraxiom}{Projective Reality}
Observable spacetime $M^{(4)}$ represents a four-dimensional projection of a higher-dimensional information space $\mathcal{I}$:
\begin{equation}
\pi: \mathcal{I} \rightarrow M^{(4)}
\label{eq:axiom1}
\end{equation}
This projection mapping $\pi$ establishes the fundamental relationship between the unobservable information substrate and measurable physical phenomena.
\end{qraxiom}

\begin{qraxiom}{Information Conservation}
The total information $I_{\text{total}}$ in the information space remains conserved, while local information distribution undergoes projective fluctuations:
\begin{equation}
\frac{dI_{\text{total}}}{dt} = 0, \quad \text{but} \quad \frac{\partial I(\vec{x},t)}{\partial t} \neq 0
\label{eq:axiom2}
\end{equation}
\end{qraxiom}

\begin{qraxiom}{Fractal Scaling}
The projection follows a fractal hierarchy with characteristic dimension $d_f$:
\begin{equation}
\mathcalO(r) \sim r^{d_f-4}, \quad \text{where} \quad d_f = 4 + \epsilon, \quad \epsilon \ll 1
\label{eq:axiom3}
\end{equation}
The small parameter $\epsilon \sim 10^{-4}$ quantifies deviations from exact four-dimensionality.
\end{qraxiom}

\begin{qraxiom}{String-Topological Corrections}
Projective dynamics incorporate topological properties of the underlying string manifold:
\begin{equation}
\mathcalO_{\text{total}} = \mathcalO_{\QR} \cdot g_s^{\chi(\mathcal{M})}
\label{eq:axiom4}
\end{equation}
where $\chi(\mathcal{M})$ denotes the Euler characteristic and $g_s$ the string coupling constant.
\end{qraxiom}

\begin{qraxiom}{Dynamic Projection}
The projection operates through differential operators acting on the coherent information density $\Inu(t)$:
\begin{equation}
\qroperator{n} = \frac{\partial^n}{\partial t^n} \log \Inu(t)
\label{eq:axiom5}
\end{equation}
\end{qraxiom}

\begin{qraxiom}{Entropic Modulation}
Projective weighting follows the local entropy density:
\begin{equation}
\mathcalW \propto \exp\left(\frac{\mathcalS}{k_B}\right)
\label{eq:axiom6}
\end{equation}
where $\mathcalS$ represents the local entropy.
\end{qraxiom}

\begin{qraxiom}{Holographic Principle}
Maximum information content is bounded by surface area:
\begin{equation}
I_{\max} = \frac{A}{4\ell_P^2}
\label{eq:axiom7}
\end{equation}
\end{qraxiom}

\begin{qraxiom}{Metric Projection}
The spacetime metric emerges from information density distribution:
\begin{equation}
g_{\mu\nu} = \eta_{\mu\nu} + \kappa \frac{\partial^2 \log \Inu}{\partial x^\mu \partial x^\nu}
\label{eq:axiom8}
\end{equation}
\end{qraxiom}

\subsection{Variational Principle and Field Dynamics}
\label{subsec:variational-principle}

The temporal evolution of the information density $\Inu(t)$ follows from a variational principle derived from the axiomatic structure. We postulate an effective action for the coupled system of scale factor $a(t)$ and information field $\phi \equiv \log \Inu$:

\begin{equation}
S[\phi, a] = \int dt \, a^3(t) \left[ \frac{\kappa_\pi}{2} \left(\frac{\partial \phi}{\partial t}\right)^2 - U(\phi) \right]
\label{eq:qr-action}
\end{equation}

The dimensionless coupling $\kappa_\pi$ normalizes to the projective length scale $\lpoi$, while $U(\phi)$ represents an effective potential determined by consistency with early-universe observations.

Variation with respect to $\phi$ yields the Euler-Lagrange equation:
\begin{equation}
\frac{d}{dt}\left(a^3 \frac{\partial \phi}{\partial t}\right) + \frac{a^3}{\kappa_\pi} \frac{dU}{d\phi} = 0
\label{eq:euler-lagrange}
\end{equation}

Substituting $\phi = \log \Inu$ and $H = \dot{a}/a$ gives:
\begin{equation}
\dot{H} + 3H^2 + \frac{1}{\kappa_\pi} U'(\phi) = 0
\label{eq:modified-friedmann}
\end{equation}

\subsection{Potential Specification and Parameter Reduction}
\label{subsec:potential-specification}

Physical consistency requires that the QR evolution reduces to standard radiation-plus-matter behavior at high redshifts $z \gg 10^3$:

\begin{equation}
H^2(a) \xrightarrow{a \ll 1} H_0^2 \left[ \Omega_r a^{-4} + \Omega_m a^{-3} \right]
\label{eq:high-z-limit}
\end{equation}

This constraint determines the potential derivative:
\begin{equation}
U'(\phi) = -\frac{3\kappa_\pi}{2} H^2 + 4\pi G \rho_{\text{eff}}(a)
\label{eq:potential-derivative}
\end{equation}

where $\rho_{\text{eff}}(a) = \rho_{r,0} a^{-4} + \rho_{m,0} a^{-3}$ represents the effective matter-radiation density.

\subsection{Information Density and Cosmological Functions}
\label{subsec:cosmological-functions}

The information density follows from the fundamental relation $H = \dot{\phi} = d\log\Inu/dt$:

\begin{align}
\phi(t) &= \phi(t_0) + \int_{t_0}^t H(t') dt' \\
\Inu(t) &= \Inu(t_0) \exp\left(\int_{t_0}^t H(t') dt'\right)
\label{eq:info-density-evolution}
\end{align}

The transformation from cosmic time $t$ to redshift $z$ uses the kinematic relation $dt/dz = -1/[(1+z)H(z)]$, enabling calculation of observable distance measures:

\begin{align}
D_C(z) &= c \int_0^z \frac{dz'}{H(z')} \\
d_L(z) &= (1+z) D_C(z)
\label{eq:distance-measures}
\end{align}

\subsection{Projective Operators and Emergent Gravity}
\label{subsec:projective-operators}

The projection from information space to observable phenomena proceeds through a hierarchy of operators that build complexity through successive applications:

\textbf{Level 1 - Fractal Enhancement:}
\begin{equation}
\qroperator{1} = \frac{\partial}{\partial t} \log \Inu(t) \cdot \mathcalF(r)
\label{eq:level1-operator}
\end{equation}

where $\mathcalF(r) = (r/\lpoi)^{d_f-4}$ provides fractal weighting based on the characteristic length scale.

\textbf{Level 2 - Spatial Projection:}
\begin{equation}
\qroperator{2} = \frac{\partial^2}{\partial t^2} \log \Inu(t) \cdot \mathbb{P}(\vec{x})
\label{eq:level2-operator}
\end{equation}

The projective function $\mathbb{P}(\vec{x})$ maps coordinate locations to observable regions of spacetime.

\textbf{Level 3 - String-Topological Modulation:}
\begin{equation}
\qroperator{3} = \qroperator{2} \cdot g_s^{\chi(\mathcal{M})} \cdot Z_{\text{String}}
\label{eq:level3-operator}
\end{equation}

This level incorporates string path integral contributions $Z_{\text{String}}$ and topological corrections.

\textbf{Level 4 - Observable Spacetime Dynamics:}
\begin{equation}
\mathcalO_{\text{total}} = \qroperator{3} \cdot \projweight{\rho_{\text{bar}}, |\psi|, r}
\label{eq:total-operator}
\end{equation}

The final weighting function $\mathcalW$ depends on baryon density, quantum state normalization, and spatial coordinates, producing the observable gravitational effects.

\subsection{Modified Einstein Equations}
\label{subsec:modified-einstein}

The complete QR field equations take the form:

\begin{equation}
\qrfieldequation
\label{eq:modified-einstein}
\end{equation}

The effective stress-energy tensor combines standard matter with QR contributions:

\begin{equation}
T_{\mu\nu}^{\text{eff}} = T_{\mu\nu}^{\text{matter}} + T_{\mu\nu}^{\QR}
\label{eq:effective-stress-energy}
\end{equation}

The QR contribution follows from the parameter-reduced form with $\beta = 1$ and $\gamma = (\log \Inu)^2$:

\begin{equation}
\qrstressenergy
\label{eq:qr-stress-energy}
\end{equation}

\subsection{Consistency Checks and Limiting Behavior}
\label{subsec:consistency-checks}

The QR equations reduce to general relativity in appropriate limits while providing specific deviations that account for observed phenomena. In the weak-field, slow-motion limit relevant to galactic dynamics, the theory produces modified Poisson equations that generate flat rotation curves without dark matter.

For cosmological applications, the QR evolution equation naturally incorporates both matter and radiation epochs while providing smooth transitions that resolve the Hubble tension through dynamic information redistribution. The fractal corrections remain small ($\epsilon \sim 10^{-4}$) but provide sufficient modification to account for observed large-scale structure statistics.

The string-topological corrections ensure compatibility with established results in theoretical physics while opening new avenues for experimental verification through gravitational wave observations and cosmic microwave background measurements.